\documentclass[12pt,a4paper]{article}
\usepackage[T2A]{fontenc}
\usepackage[utf8]{inputenc}
\usepackage[russian]{babel}
\usepackage{amsmath,amssymb}
\usepackage{geometry}
\geometry{margin=2.5cm}

\begin{document}

\begin{center}
{\Large\textbf{Алгоритм анализа пространственного распределения шипиков вдоль полной дендритной сети нейрона}}\\[0.5em]
{\normalsize Теоретическое и вычислительное описание конвейера, основанного на 3D network spatial analysis}
\end{center}

\noindent\textbf{Постановка задачи.}
Цель алгоритма анализа всей сети нейрона состоит в том, чтобы описать не отдельный дендритный фрагмент, а распределение дендритных шипиков вдоль полной дендритной арборизации или вдоль крупной части арборизации, представленной как связная трёхмерная сеть. В отличие от локального анализа одного дендрита, где основной объект исследования есть набор шипиков на поверхности одного меша, сетевой анализ рассматривает шипики как события точечного процесса, ограниченного дендритным деревом. Такая постановка следует работе Anton-Sanchez et al. ``Three-dimensional spatial modeling of spines along dendritic networks in human cortical pyramidal neurons'', где подчёркивается, что шипики не расположены свободно в трёхмерном пространстве, а могут лежать только на дендритном стволе. Поэтому прямые евклидовы расстояния между шипиками не являются корректной основной метрикой: две точки могут быть близки в \(\mathbb{R}^3\), но находиться на разных ветвях и быть далеко друг от друга по дендритной сети.

\noindent\textbf{Линейная сеть как математическая модель дендрита.}
Полная дендритная арборизация моделируется как линейная сеть \(L\), то есть как объединение конечного числа линейных сегментов:
\[
L=\bigcup_{e=1}^{m} l_e.
\]
Каждый сегмент \(l_e=[u_e,v_e]\) задаётся двумя концами \(u_e,v_e\in\mathbb{R}^3\) и представляет часть центральной оси дендрита. Общая длина сети равна
\[
|L|=\sum_{e=1}^{m}|l_e|.
\]
Если сеть связна и не содержит циклов, она является деревом, что соответствует типичной топологии дендритной арборизации. Основной геометрической метрикой является кратчайшее расстояние по сети
\[
d_L(u,v)=\min_{\pi:u\rightarrow v}\int_{\pi} ds,
\]
где минимум берётся по всем путям \(\pi\), соединяющим точки \(u\) и \(v\) внутри сети \(L\). Если путь между точками отсутствует, расстояние считается бесконечным. Именно \(d_L\), а не евклидова норма \(\|u-v\|_2\), используется для анализа близости шипиков, интенсивности точечного процесса и функции Рипли.

\noindent\textbf{Шипики как события точечного процесса на сети.}
Набор шипиков представляется как реализация точечного процесса
\[
X=\{x_1,\ldots,x_n\}, \qquad x_i\in L,
\]
где \(x_i\) есть точка крепления \(i\)-го шипика к дендритному стволу, спроецированная на центральную ось сети. В биологической интерпретации \(x_i\) задаёт место потенциального возбуждающего синаптического входа на дендритной арборизации. В статистической интерпретации это событие на линейной сети. Такой переход позволяет анализировать не только абсолютное число шипиков, но и то, как они распределены вдоль ветвления, как меняется их плотность при удалении от сомы и отличаются ли правила распределения между апикальными и базальными деревьями или между разными нейронами.

\noindent\textbf{Соответствие алгоритма статье Anton-Sanchez et al.}
В статье исходными объектами являлись трёхмерные реконструкции пирамидных нейронов, из которых извлекались оси дендритных деревьев и точки крепления шипиков. Далее применялся анализ точечных паттернов на трёхмерных линейных сетях: оценивалась интенсивность шипиков вдоль сети, проверялась зависимость интенсивности от расстояния до тела клетки, подгонялась неоднородная пуассоновская модель с лог-квадратичной зависимостью от расстояния до сомы, вычислялась геометрически скорректированная неоднородная сетевая функция Рипли \(K\), строились Monte Carlo envelopes и выполнялось сравнение групп с помощью replicated point patterns. В текущем проекте основные строительные блоки этого подхода реализованы в файле \texttt{dendrite\_analysis/network.py}: построение дендритного графа, проекция шипиков на граф, сетевые расстояния, оценка интенсивности, неоднородная пуассоновская модель, network Ripley \(K\), симуляции пуассоновского процесса на графе, envelopes, сравнение групп и генерация отчёта. Однако этот модуль пока является самостоятельным конвейером и не встроен в основной ноутбук \texttt{dendrite-analysis.ipynb}; кроме того, некоторые элементы статьи реализованы приближённо, что отдельно описано ниже.

\noindent\textbf{Построение дендритной сети из мешей.}
В вычислительной реализации дендритная сеть представлена объектом \texttt{DendriticGraph}. Это взвешенный неориентированный граф \(G=(V,E)\), где вершины соответствуют узлам скелета дендрита, а рёбра соответствуют сегментам центральной оси. Каждая вершина хранит трёхмерную координату, тип узла и идентификатор дендрита, а каждое ребро хранит длину. Если на вход передаётся один меш дендрита, функция \texttt{build\_dendritic\_graph} выполняет скелетизацию меша и переводит полученные скелетные сегменты в граф. Если на вход подаётся несколько дендритных мешей, функция \texttt{build\_graph\_from\_meshes} строит подграфы для каждого дендрита и объединяет их в одну сеть; близкие концевые узлы разных подграфов могут сшиваться при помощи порога \texttt{snap\_threshold}. Таким образом, на уровне низкоуровневых функций уже предусмотрена работа не только с одним дендритом, но и с сетью, состоящей из нескольких дендритных деревьев или ветвей.

\noindent\textbf{Сома, корень дерева и расстояние до тела клетки.}
Для анализа неоднородности распределения шипиков выбирается корень сети \(r\), соответствующий соме или ближайшему к соме узлу дендритного графа. Если координата сомы передана явно, ближайший узел графа помечается как \texttt{soma}; если корень не задан, код использует первый узел графа как fallback и выдаёт предупреждение. После выбора корня для всех узлов вычисляются расстояния
\[
d_L(u,r),
\]
то есть длины кратчайших путей от сомы до узлов сети. Для произвольной точки внутри ребра расстояние до сомы вычисляется через расстояния до двух концов ребра и положение точки на этом ребре. Эта величина является основным пространственным ковариатом в модели интенсивности, поскольку в биологическом смысле плотность шипиков часто меняется вдоль проксимально-дистальной оси дендритного дерева.

\noindent\textbf{Проекция шипиков на дендритный граф.}
Исходные точки крепления шипиков обычно заданы в трёхмерных координатах реконструкции. Чтобы использовать методы point process on network, каждая точка \(p_i\in\mathbb{R}^3\) проецируется на ближайшее ребро графа. Для ускорения поиска строится KD-дерево по серединам рёбер, после чего для каждого шипика рассматривается несколько ближайших кандидатных рёбер. Для ребра \([u,v]\) вычисляется ортогональная проекция
\[
\hat p_i = u + t(v-u), \qquad
t=\frac{(p_i-u)^\top(v-u)}{\|v-u\|^2}, \qquad t\in[0,1].
\]
Выбирается ребро, дающее минимальное расстояние \(\|p_i-\hat p_i\|_2\). Если это расстояние не превышает \texttt{max\_distance\_to\_edge}, шипик считается успешно сопоставленным сети; иначе он попадает в список неприсвоенных объектов. Для каждого спроецированного шипика сохраняются исходная точка, проекция на сеть, идентификатор ребра, параметр положения \(t\), расстояние до ребра и расстояние до сомы.

\noindent\textbf{Матрица сетевых расстояний между шипиками.}
После проекции шипиков вычисляется матрица попарных расстояний \(D=(d_{ij})\), где
\[
d_{ij}=d_L(x_i,x_j).
\]
Если два шипика лежат на одном ребре, расстояние между ними равно длине соответствующего участка ребра. Если они лежат на разных рёбрах, рассматриваются четыре варианта пути через концы этих рёбер, и выбирается минимальный:
\[
d_{ij}=\min
\begin{cases}
t_iL_i+d_L(u_i,u_j)+t_jL_j,\\
t_iL_i+d_L(u_i,v_j)+(1-t_j)L_j,\\
(1-t_i)L_i+d_L(v_i,u_j)+t_jL_j,\\
(1-t_i)L_i+d_L(v_i,v_j)+(1-t_j)L_j.
\end{cases}
\]
Здесь \(u_i,v_i\) есть концы ребра, на котором лежит \(x_i\), \(L_i\) длина этого ребра, а \(t_i\) положение точки на ребре. Расстояния между узлами графа вычисляются алгоритмом Дейкстры. Эта матрица является общей основой для функции Рипли, групповых сравнений и любых дополнительных статистик взаимного расположения шипиков на уровне всей сети.

\noindent\textbf{Глобальная плотность и однородная CSR-модель.}
Самая простая null-модель предполагает однородный пуассоновский процесс на сети, или Complete Spatial Randomness on a network. В этом случае интенсивность постоянна:
\[
\lambda(u)\equiv \lambda,
\qquad
\hat\lambda=\frac{n}{|L|}.
\]
Здесь \(n\) есть число шипиков, а \(|L|\) общая длина дендритной сети. Эта величина интерпретируется как среднее число шипиков на единицу длины дендритной оси. Однородная модель удобна как базовый ориентир, но для дендритных шипиков она часто биологически недостаточна, поскольку плотность шипиков обычно зависит от расстояния до тела клетки. Именно поэтому в статье сначала анализируется пространственная неоднородность интенсивности, а затем уже проверяется кластеризация относительно неоднородной модели.

\noindent\textbf{Бинированная оценка интенсивности вдоль проксимально-дистальной оси.}
Функция \texttt{estimate\_binned\_intensity} оценивает зависимость плотности шипиков от расстояния до сомы. Ось расстояний разбивается на интервалы \([a_k,b_k)\). Для каждого интервала считается число шипиков
\[
N_k=\#\{i:a_k\leq d_L(x_i,r)<b_k\}
\]
и длина части сети, попадающая в тот же диапазон расстояний до сомы. В реализации эта длина оценивается численно: сеть дискретизируется набором точек, для них вычисляется расстояние до сомы, а затем доля точек в бине умножается на общую длину сети. Интенсивность в бине равна
\[
\hat\lambda_k=\frac{N_k}{|L_k|},
\]
где \(|L_k|\) есть длина сети в соответствующем диапазоне расстояний. Такая оценка даёт грубую, но наглядную картину того, где вдоль арборизации шипики расположены плотнее, а где разреженнее.

\noindent\textbf{Сглаженная оценка интенсивности.}
Функция \texttt{estimate\_smooth\_intensity} строит сглаженную оценку \(\lambda(d)\), где \(d=d_L(u,r)\). Идея состоит в том, что наблюдаемое распределение расстояний шипиков до сомы нужно нормировать на доступную длину сети при каждом расстоянии. Если \(\hat f_X(d)\) есть ядерная оценка плотности расстояний шипиков до сомы, а \(\hat f_L(d)\) описывает, сколько дендритной длины доступно на расстоянии \(d\), то интенсивность приближённо оценивается как
\[
\hat\lambda(d)=\frac{\hat f_X(d)}{\hat f_L(d)}.
\]
В коде \(\hat f_X(d)\) вычисляется гауссовым ядром, а \(\hat f_L(d)\) аппроксимируется по дискретизации сети. Такая нормировка важна, потому что большое число шипиков на некотором расстоянии может быть связано не только с высокой плотностью, но и с тем, что на этом расстоянии находится большая суммарная длина ветвей.

\noindent\textbf{Проверка зависимости интенсивности от расстояния до сомы.}
В статье Anton-Sanchez et al. для проверки зависимости интенсивности от расстояния до тела клетки используется CDF-тест: распределение значений ковариаты \(d_L(u,r)\) в точках шипиков сравнивается с распределением этой же ковариаты у случайных точек на сети. В текущей реализации эта идея представлена функцией \texttt{test\_intensity\_dependence\_cdf}: сеть равномерно дискретизируется вдоль рёбер, после чего распределение расстояний до сомы у шипиков сравнивается с распределением расстояний до сомы у точек сети с помощью двухвыборочного критерия Колмогорова--Смирнова. Если эти распределения статистически различаются, интенсивность шипиков не может считаться однородной относительно расстояния до сомы. Дополнительно сохраняется likelihood-ratio диагностика \texttt{test\_intensity\_dependence}, в которой после бинирования подгоняются две пуассоновские модели для чисел шипиков в бинах. Нулевая модель предполагает постоянную интенсивность, а альтернативная модель добавляет линейную зависимость от расстояния:
\[
\log \mu_k = \log |L_k|+\theta_0
\]
для постоянной модели и
\[
\log \mu_k = \log |L_k|+\theta_0+\theta_1 d_k
\]
для модели с расстоянием, где \(d_k\) есть середина \(k\)-го бина, а \(\mu_k\) ожидаемое число шипиков в бине. Далее вычисляется likelihood-ratio statistic
\[
LR=2(\ell_1-\ell_0),
\]
и \(p\)-значение получается по распределению \(\chi^2\) с одной степенью свободы. Таким образом, CDF/KS-тест является более близким к референсной статье этапом, а likelihood-ratio тест остаётся дополнительной модельной проверкой направления и силы зависимости.

\noindent\textbf{Неоднородная пуассоновская модель.}
После обнаружения неоднородности интенсивности подгоняется неоднородный пуассоновский процесс на сети. В реализации функция \texttt{fit\_inhomogeneous\_poisson} использует лог-линейную модель
\[
\lambda(u)=\exp(\theta^\top z(u)),
\]
где \(z(u)\) есть вектор ковариат точки сети. По умолчанию используются константа, расстояние до сомы и квадрат расстояния до сомы:
\[
z(u)=\left(1,d_L(u,r),d_L(u,r)^2\right).
\]
Тем самым получается лог-квадратичная модель, совпадающая по идее с моделью статьи:
\[
\lambda(u)=\exp\left(\theta_0+\theta_1d_L(u,r)+\theta_2d_L(u,r)^2\right).
\]
Лог-правдоподобие пуассоновского точечного процесса на сети имеет вид
\[
\ell(\theta)=\sum_{i=1}^{n}\log\lambda(x_i)-\int_L \lambda(u)\,du
=
\sum_{i=1}^{n}\theta^\top z(x_i)-\int_L\exp(\theta^\top z(u))\,du.
\]
Интеграл по сети в коде оценивается численно по точкам, равномерно сэмплированным вдоль рёбер графа. После оптимизации сохраняются коэффициенты, стандартные ошибки, лог-правдоподобие, AIC, BIC, значения подогнанной интенсивности в точках шипиков и остатки. Интерпретационно эта модель отделяет неоднородность, связанную с проксимально-дистальным положением, от дополнительной пространственной агрегации шипиков.

\noindent\textbf{Сетевая функция Рипли.}
Функция Рипли \(K\) является сводной характеристикой второго порядка. В классической пространственной статистике она измеряет ожидаемое число других точек в радиусе \(r\) вокруг типичной точки, нормированное на интенсивность. Для сети евклидово расстояние заменяется кратчайшим расстоянием \(d_L\). В геометрически скорректированном варианте текущая реализация вычисляет multiplicity-функцию \(m(u,t)\), то есть число точек сети, находящихся от \(u\) ровно на сетевом расстоянии \(t\). Тогда однородная оценка имеет вид
\[
\hat K_L(r)=\frac{|L|}{n^2}\sum_{i=1}^{n}\sum_{j=1}^{n}
\frac{\mathbf{1}\{d_L(x_i,x_j)\leq r,\ i\neq j\}}{m(x_i,d_L(x_i,x_j))}.
\]
Если \(\hat K_L(r)\) выше ожидаемой кривой, точки расположены ближе друг к другу, чем ожидалось от случайного процесса, что интерпретируется как кластеризация. Если \(\hat K_L(r)\) ниже ожидаемой кривой, это соответствует регулярности или отталкиванию. При homogeneous CSR для геометрически скорректированной сетевой функции Рипли ожидаемая кривая равна
\[
K_{\mathrm{exp}}(r)=r.
\]
В коде \(m(u,t)\) вычисляется на metric graph: для каждого ребра рассматривается нижняя огибающая двух линейных функций расстояния через концы ребра, а затем подсчитываются пересечения этой функции с уровнем \(t\). Это делает параметр \texttt{correction="geometric"} содержательной геометрической поправкой, а не глобальным множителем на доступную длину.

\noindent\textbf{Неоднородная сетевая функция Рипли.}
Если интенсивность вдоль сети непостоянна, простое превышение \(K(r)\) над ожидаемой кривой может отражать не кластеризацию, а изменение плотности с расстоянием до сомы. Поэтому используется неоднородная версия функции Рипли, в которой каждая пара точек взвешивается обратным произведением интенсивностей и геометрической multiplicity-поправкой:
\[
\hat K_{\mathrm{inhom}}(r)=
\frac{1}{|L|}
\sum_{i=1}^{n}\sum_{j=1}^{n}
\frac{\mathbf{1}\{d_L(x_i,x_j)\leq r,\ i\neq j\}}
{\hat\lambda(x_i)\hat\lambda(x_j)m(x_i,d_L(x_i,x_j))}.
\]
Здесь \(\hat\lambda(x_i)\) берётся из подогнанной пуассоновской модели. Эта статистика проверяет, остаётся ли пространственная структура после того, как учтена проксимально-дистальная неоднородность. В биологической интерпретации это позволяет различить два сценария: шипики могут казаться сгруппированными просто потому, что на определённом расстоянии от сомы их в целом больше, либо они действительно образуют дополнительные локальные скопления относительно ожидаемой плотности.

\noindent\textbf{Симуляция пуассоновского процесса на дендритной сети.}
Для построения Monte Carlo envelopes код использует функцию \texttt{simulate\_poisson\_on\_graph}. В однородном случае рёбра выбираются с вероятностями, пропорциональными их длинам, после чего положение точки на выбранном ребре сэмплируется равномерно. Это даёт равномерный процесс по длине сети. В неоднородном случае используется acceptance-rejection схема: сначала генерируются кандидаты равномерно по сети, затем кандидат принимается с вероятностью, пропорциональной \(\lambda(d_L(u,r))\). Число симулированных точек фиксируется равным наблюдаемому числу шипиков \(n\), что соответствует условной симуляции при заданной мощности паттерна.

\noindent\textbf{Monte Carlo envelopes и проверка отклонения от модели.}
Функция \texttt{compute\_simulation\_envelopes} строит наблюдаемую \(K\)-кривую, затем многократно симулирует точечные паттерны из однородной или неоднородной пуассоновской модели и вычисляет \(K\)-кривую для каждой симуляции. Для каждого значения \(r\) берутся нижний и верхний квантили симуляционного распределения. Дополнительно вычисляется глобальная статистика отклонения
\[
T_{\mathrm{obs}}=\max_r |\hat K_{\mathrm{obs}}(r)-K_{\mathrm{exp}}(r)|
\]
и аналогичные значения \(T_b\) для симуляций. \(p\)-значение оценивается как доля симуляций, у которых \(T_b\geq T_{\mathrm{obs}}\). Если наблюдаемая кривая существенно выше envelope, распределение интерпретируется как кластеризованное относительно выбранной null-модели. Если она существенно ниже, это указывает на регулярность. Если кривая остаётся внутри envelope, null-модель не отвергается. В статье использовались global envelopes на основе 19 симуляций для уровня \(0.05\); в коде число симуляций задаётся параметром \texttt{n\_simulations}, по умолчанию равным \(99\) в высокоуровневом pipeline.

\noindent\textbf{Сравнение групп и replicated point patterns.}
В статье дендритные деревья рассматриваются как реплики точечного процесса: например, несколько базальных деревьев одного нейрона являются несколькими реализациями одного типа распределения. Это позволяет сравнивать группы паттернов, а не только отдельные деревья. В текущем коде функция \texttt{group\_k\_analysis} вычисляет \(K\)-кривые для набора групп, а \texttt{permutation\_test\_groups} сравнивает две группы по интегральной разности средних \(K\)-кривых:
\[
T=\int_{r_{\min}}^{r_{\max}}
\left|\bar K_A(r)-\bar K_B(r)\right|\,dr.
\]
В дискретной реализации интеграл заменяется суммой по сетке \(r\). Затем паттерны перемешиваются между группами, для каждой перестановки пересчитывается статистика \(T\), и \(p\)-значение определяется как доля перестановок с отклонением не меньше наблюдаемого. Эта процедура близка по смыслу к анализу replicated point patterns, но не является точной копией studentized permutation test из статьи, поскольку текущая статистика не использует полную studentized-нормировку внутригрупповых вариаций \(K\)-кривых.

\noindent\textbf{Дополнительные сетевые расстояния до ветвлений и терминалей.}
Помимо попарных расстояний между шипиками модуль содержит функции \texttt{compute\_distance\_to\_branching} и \texttt{compute\_distance\_to\_terminals}. Первая вычисляет для каждого шипика минимальное сетевое расстояние до ближайшей точки ветвления или сомы, вторая до ближайшего терминального узла. Эти величины важны для биологической интерпретации, поскольку локальное положение шипика относительно ветвлений может влиять на интеграцию синаптических входов, а расстояние до терминального конца характеризует проксимально-дистальную организацию внутри конкретной ветви. В текущем высокоуровневом \texttt{run\_analysis} эти функции не являются центральным этапом отчёта, но они уже доступны как готовые компоненты для расширения нейронного анализа.

\noindent\textbf{Визуализация и сохраняемые результаты.}
Модуль \texttt{network.py} содержит функции визуализации трёхмерной сети, интенсивности по расстоянию до сомы, \(K\)-кривой, отклонения \(K(r)-K_{\mathrm{exp}}(r)\) и группового сравнения \(K\)-кривых. Высокоуровневая функция \texttt{generate\_report} собирает краткую статистику числа шипиков, общей длины сети, плотности, среднего расстояния до сомы, числа узлов и рёбер, параметры пуассоновской модели и результат \(K\)-анализа. Если доступны графические библиотеки, дополнительно сохраняются графики интенсивности, функции Рипли, отклонения от ожидания и интерактивная 3D-визуализация дендритной сети со шипиками.

\noindent\textbf{Высокоуровневый pipeline \texttt{run\_analysis}.}
Функция \texttt{run\_analysis} является автономным входом в сетевой анализ одного переданного дендритного объекта. Она принимает путь к мешу или уже загруженный меш дендрита, словарь точек крепления шипиков, опциональную координату сомы, путь к конфигурации и папку вывода. Далее она строит граф дендрита, проецирует шипики на граф, оценивает бинированную и сглаженную интенсивность, подгоняет неоднородную пуассоновскую модель, строит сетевую \(K\)-функцию с симуляционными envelopes и генерирует HTML- или Markdown-отчёт. Для анализа полной MICrONS-структуры добавлен объектный слой \texttt{dendrite\_analysis/neuron.py}: классы \texttt{Soma}, \texttt{Limb}, \texttt{Branch} и \texttt{Neuron} читают иерархию нейрона, строят общий граф из \texttt{limb\_skeleton.npy} или \texttt{branch\_skeleton.npy}, собирают точки крепления шипиков из папок \texttt{spines} и запускают тот же сетевой анализ уже на уровне всего нейрона. На этом же объектном слое теперь доступен единый режим запуска \texttt{Neuron.run\_full\_analysis(mode)}, где \texttt{mode="network"} выполняет только сетевой анализ нейрона, \texttt{mode="branches"} выполняет только анализ отдельных дендритных branch-фрагментов через класс \texttt{DendriteBranch}, а \texttt{mode="both"} последовательно запускает оба уровня анализа на уже загруженных мешах и скелетах без повторного чтения данных с диска. Для пакетной обработки корневой папки MICRONS добавлена функция \texttt{run\_microns\_neuron\_analyses}.

\noindent\textbf{Текущий статус реализации в проекте.}
На уровне локального кода модуль \texttt{dendrite\_analysis/network.py} содержит реализацию алгоритма, вдохновлённого статьёй Anton-Sanchez et al., а модуль \texttt{dendrite\_analysis/neuron.py} добавляет интеграционный слой для MICRONS-нейронов. Для запуска сетевого, branch-level или объединённого анализа создан отдельный ноутбук \texttt{neuron-network-analysis.ipynb}; режим выбирается переменной \texttt{ANALYSIS\_MODE}. Итоговый файл \texttt{neuron\_structural\_network\_vectors.csv} содержит одну строку на нейрон и хранит core ML-вектор: структурные размеры сети, missing-mask признаки, общую и типоспецифичную линейную плотность шипиков, ветвлений и терминальных узлов для апикальных и базальных частей, характеристики сомы, effect-size статистики зависимости интенсивности от расстояния до сомы, нормированные коэффициенты неоднородной пуассоновской модели и скалярные effect-size агрегаты network \(K\)-анализа. \(p\)-значения статистических тестов, параметры Monte Carlo envelope, качество проекции и технические счётчики сохраняются отдельно в \texttt{neuron\_structural\_network\_metadata.csv}; raw log-likelihood, AIC/BIC, их нормированные версии и стандартные ошибки коэффициентов сохраняются в \texttt{neuron\_structural\_network\_fit\_quality.csv}. Network \(K\)-анализ выполняется как для всей сети нейрона, так и отдельно для апикальной и базальной частей, если в соответствующем compartment есть как минимум три спроецированных шипика и ненулевая длина сети; в ML-вектор попадают только скалярные агрегаты этих кривых. Полноценное replicated point pattern comparison между апикальной и базальной группами пока не выполняется, поскольку текущая цель состоит не в проверке групповой гипотезы, а в построении подробного признакового описания нейрона. Полные кривые network \(K\) сохраняются для каждого нейрона в \texttt{neuron\_structural\_network\_ripley\_k\_network.csv}, \texttt{neuron\_structural\_network\_ripley\_k\_network\_apical.csv} и \texttt{neuron\_structural\_network\_ripley\_k\_network\_basal.csv}, если соответствующий анализ был выполнен, а не разворачиваются в итоговый ML-вектор. При запуске branch-level части из того же объекта \texttt{Neuron} дополнительно формируются файлы локального анализа дендритных фрагментов, включая \texttt{dendrite\_structural\_organization\_vectors.csv} и \texttt{all\_dendr\_metrics.csv} в отдельной папке вывода. Старый branch-level pipeline \texttt{dendrite-analysis.ipynb} остаётся самостоятельным способом запуска анализа отдельных дендритов; класс \texttt{Dendrite} сохранён для обратной совместимости, а имя \texttt{DendriteBranch} добавлено как семантический alias.

\noindent\textbf{Отличия текущей реализации от статьи.}
Первое отличие состоит в способе получения сети. В статье оси дендритов и точки крепления извлекались из реконструкций в формате \texttt{.vrml}; в проекте сеть строится из готовых \texttt{.npy}-скелетов MICRONS, из CGAL-скелетизации мешей или из уже переданной графовой структуры. Второе отличие состоит в технической реализации CDF-проверки: статья формулирует CDF-тест в терминах ковариат на линейной сети, а текущий код реализует близкий по смыслу двухвыборочный KS-тест между расстояниями шипиков до сомы и расстояниями равномерно сэмплированных точек сети до сомы. Третье отличие состоит в групповом сравнении: статья использует studentized permutation test для replicated point patterns, тогда как текущая функция \texttt{permutation\_test\_groups} сравнивает интегральную разность средних \(K\)-кривых без полной studentized-стандартизации. Основное ядро, влияющее на выводы о кластеризации, а именно неоднородная модель интенсивности, геометрически скорректированная network \(K_L\)-функция и Monte Carlo envelopes, теперь реализовано в референсно близком виде.

\noindent\textbf{Интерпретация результата сетевого анализа.}
Результат анализа всей сети нейрона должен интерпретироваться на нескольких уровнях. Глобальная плотность \(n/|L|\) показывает среднее число шипиков на единицу длины дендритной арборизации, а типоспецифичные плотности отдельно описывают апикальные и базальные части сети. Бинированная и сглаженная интенсивность показывают, как эта плотность меняется при удалении от сомы. Неоднородная пуассоновская модель проверяет, можно ли описать распределение шипиков как случайное размещение на сети с интенсивностью, зависящей от расстояния до тела клетки. Network Ripley \(K\) и Monte Carlo envelopes показывают, остаётся ли дополнительная кластеризация или регулярность после учёта такой неоднородности. Групповые перестановочные сравнения позволяют проверять, отличаются ли правила пространственного распределения между апикальными и базальными деревьями, между нейронами или между экспериментальными группами. В совокупности эти этапы переводят реконструкцию нейрона в статистическое описание его сетевой организации.

\noindent\textbf{Связь с анализом отдельных дендритов.}
Локальный анализ, описанный в отчёте \texttt{dendrite\_spatial\_morphology\_analysis\_report.tex}, и сетевой анализ решают разные, но взаимодополняющие задачи. Локальный анализ работает на уровне одного дендритного фрагмента и строится вокруг расстояний между точками крепления по поверхности меша, DBSCAN-кластеров, nearest-neighbor статистик и связи локального окружения с морфологией шипиков. Сетевой анализ работает на уровне полной арборизации и делает акцент на центральной оси дендрита, расстоянии до сомы, интенсивности точечного процесса, неоднородной пуассоновской модели и сетевой функции Рипли. В идеальном итоговом pipeline эти два уровня должны сосуществовать: сетевой анализ описывает правила размещения шипиков вдоль всего нейрона, а локальный анализ уточняет морфологическую и кластерную организацию внутри отдельных ветвей или фрагментов.

\noindent\textbf{Практический следующий шаг для интеграции.}
Практический следующий шаг состоит в проверке нового сетевого слоя на реальных MICRONS-нейронах и в уточнении метаданных апикальных и базальных лимбов. Если такая разметка будет поставляться вместе с данными, её следует передавать в \texttt{DENDRITE\_TYPE\_MAP}; если нет, автоматическая эвристика может использоваться только как предварительная подсказка, но финальная биологическая разметка должна оставаться ручной или таблично подтверждённой. После первичного прогона также имеет смысл оценить вычислительную стоимость геометрически скорректированной \(K_L\)-функции на больших нейронах и при необходимости добавить ускоренную аппроксимацию multiplicity-функции для exploratory-прогонов, сохранив точный вариант для финального статистического анализа.

\noindent\textbf{Источники.}
Основная методологическая основа отчёта: Anton-Sanchez L., Larrañaga P., Benavides-Piccione R., Fernaud-Espinosa I., DeFelipe J., Bielza C. Three-dimensional spatial modeling of spines along dendritic networks in human cortical pyramidal neurons. \textit{PLoS ONE}. 2017;12(6):e0180400. DOI: \texttt{10.1371/journal.pone.0180400}. Публичная страница статьи доступна по адресу \texttt{https://journals.plos.org/plosone/article?id=10.1371/journal.pone.0180400}, запись PubMed доступна по адресу \texttt{https://pubmed.ncbi.nlm.nih.gov/28662210/}.

\end{document}
